%表紙
\newpage
\chapter{学生}
\begin{enumerate}[label=第\arabic*条　]
  \item この章は、憲章の趣旨に基づき、学生の資格、入会、退会その他の学生の在籍に関し必要な事項を定めるものとする。
  \item この規則を遵守できる者は、学生となることができる。
    \begin{enumerate}[label=\arabic*　,start=2,leftmargin=*]
      \item 学園が活動の基盤として利用する外部サービスの定める利用条件を満たさない者は、当該サービスにおける活動に参加することができないため、前項の規定にかかわらず、学生となることができない。
    \end{enumerate}
  \item 学生となろうとする者は、所定の申請を行うものとし、学園が新規申請を受け付けている間に当該申請が学園に到達したときに、学生としての資格を取得する。申請について、認可又は選考を要しない。
    \begin{enumerate}[label=\arabic*　,start=2,leftmargin=*]
      \item 教務主事は、運営上必要があるときは、第３章第３条の募集期間ごとに新たに学生資格を取得する者の人数の上限を定めることができる。
      \item 前項の上限に達したときは、教務主事は、その後の新規申請の受付を終了し、終了の時点を速やかに公示する。
      \item 前二項の上限は、当該募集期間に新たに学生となる者に限るものとし、現に学生である者の在籍の継続又は総在籍者数を制限しない。
      \item 第１項の申請の具体的な手続については、別に定める。
    \end{enumerate}
  \item 学生は、憲章の定めるところにより、学園における探究活動に自由に参加することができる。
    \begin{enumerate}[label=\arabic*　,start=2,leftmargin=*]
      \item 学生は、この規則を遵守し、コミュニティの秩序の維持に努めなければならない。
    \end{enumerate}
  \item 学生は、一度入会の手続を経た後は、学期をまたいで在籍を継続するものとし、学期ごとの在籍手続を要しない。
    \begin{enumerate}[label=\arabic*　,start=2,leftmargin=*]
      \item 前項の規定は、次条の退会及び第７条の除名その他の在籍資格の喪失事由を妨げない。
    \end{enumerate}
  \item 学生は、いつでも退会することができる。
  \item 学生は、第８章に定めるところにより除名の処分を受けたときは、学生としての資格を失う。
  \item この章の実施に関し必要な事項は、教務主事が別に定める。
\end{enumerate}
